1 

Визначити, чи є одне розбиття підрозбиттям іншого. Розбиття задані як множини множин.
\\\par
2 


Для множин з повтореннями, заданими словниками, реалізувати теоретико-множинну операцію об’єднання (елемент входить у результат максимум входжень в аргументи). Розв’язання оформити у вигляді функції.
\\\par
3 





Відношення $C$ між множинами $A$ та $B$ ($A, B\subset Z$) задане словником, в якому ключами є елементи множини $A$, а значенням кожного ключа $a$ --- його повний образ $C\left(a\right)=\left\{b|\left(a,b\right)\in C\right\}$. Перевірити чи є відношення $C$ функціональним.
\\\par
4 

За розбиттям множини $A=\overline{0,n-1}$ та відображенням\\ $f{:}A\rightarrow{bool}$  побудувати підрозбиття таке, що два елементи знаходяться в одному класі підрозбиття тоді і тільки тоді, коли відображення $f$ приймає на них однакове значення. Вважати, що $f$ задано функцією. Для розбиттів використовувати зображення у вигляді множини множин.
\\\par
5 














За функціями $f$ та $g$, заданими словниками знайти їх композицію. Розв’язання оформити у вигляді функції, що приймає словники $f$ та $g$ як аргументи та результат повертає у вигляді словника.
\\\par
6 














Написати функцію, що за послідовністю цілих знаходить 5 елементів, що зустрічаються в ній найчастіше. Якщо п’яте місце розділяють кілька значень, додавати до результату всіх їх.
\\\par
7 














Бінарне відношення R задане на множині $\overline{0,n-1}$ матрицею відношення з цілими елементами.  Перевірити, чи є відношення симетричним. Обов'язково зазначити, яке зображення матриці використовується в розв'язанні: двовимірний масив типу \textbf{ndarray} або список списків.
\\\par
8 














Написати функцію, що стискає послідовність цілих, зображену списком: нулі в кінці та на початку вилучаються, у середині послідовності кожна послідовність нулів замінюється на один нуль. Наприклад, з послідовності 0,0,1,2,0,0,0,0,3,0,5,0,0 має стати 1,2,0,3,0,5. Нові об'єкти, що зберігають послідовності утворювати не можна.
\\\par


9 














Відношення $C$ між множинами $A$ та $B$ ($A, B\subset Z$) задане словником, в якому ключами є елементи множини $A$, а значенням кожного ключа $a$ --- його повний образ $C\left(a\right)=\left\{b|\left(a,b\right)\in C\right\}$. Перевірити чи є відношення $C$ ін'єктивним.
\\\par
10 














Бінарне відношення R задане на множині $\overline{0,n-1}$ матрицею відношення з цілими елементами.  Перевірити, чи є відношення антирефлексивним. Обов'язково зазначити, яке зображення матриці використовується в розв'язанні: двовимірний масив типу \textbf{ndarray} або список списків.
\\\par


11 














Відповідність $C$ між множинами $A$ та $B$ ($A\subset Z$) задана списком пар. Написати функцію, що зображує $C$ у вигляді словника, де значеннями ключів є їх повні образи.
\\\par
12 














Рядок містить слова, відокремлені крапками та комами. Знайти кількість слів, довжина яких дорівнює 13. Копій частини рядка довших 1 символу створювати не можна.
\\\par
13 














Написати функцію, що за послідовністю цілих знаходить елементи, що зустрічаються в ній рідше за інші (але зустрічаються).
\\\par
14 














Рядок містить слова, відокремлені крапками та комами. Знайти кількість слів, що містять кожну з літер s, e, l, f. Копій частини рядка довших 1 символу створювати не можна.
\\\par
\newpage
15 














Список містить цілі числа. Вилучити зі списку усі повторення (кожне значення має залишитися в одному екземплярі). Для кожного числа має залишатись його перше входження в список.
\\\par
16 














Функція $f$ між множинами $A$ та $B$ (елементи множини $A$ є такими, що хешуються) задана списком пар. Написати функцію, що зображує $f$ у вигляді словника, де значеннями ключів є значення функції.
\\\par
17 














Для множин з повтореннями, заданими словниками, реалізувати теоретико-множинну операцію перетину (елемент входить в результат мінімум входжень в аргументи). Розв’язання оформити у вигляді функції.
\\\par
18 














Бінарні відношення $R_1$, $R_2$ задано на множині $\overline{0,n-1}$ матрицями відношення з елементами логічного типу. Обчислити композицію $R_1\circ R_2$ відношень. Обов'язково зазначити, яке зображення матриці використовується в розв'язанні: двовимірний масив типу \textbf{ndarray} або список списків.
\\\par
19 














Бінарні відношення $R_1$, $R_2$ задано на множині $\overline{0,n-1}$ матрицями відношення з елементами логічного типу. Перевірити, чи правда, що $R_1\subset R_2$. Обов'язково зазначити, яке зображення матриці використовується в розв'язанні: двовимірний масив типу \textbf{ndarray} або список списків.
\\\par
20 














Побудувати функцію, що зображення розбиття множини у вигляді множини множин перетворює на зображення розбиття у вигляді словника, що задає належність елементів класам.
\\\par
21 














Бінарне відношення на множині $\overline{0,n-1}$ задане множиною пар. Перевірити, чи є відношення симетричним.
\\\par
22 














Рядок містить слова, відокремлені крапками та комами. Знайти останнє слово довжини 13. Копії частини рядка довші 1 символу можна створювати тільки в інструкції \textbf{return}.
\\\par
23 














Бінарне відношення на множині $\overline{0,n-1}$ задане множиною пар. Перевірити, чи є відношення антисиметричним.
\\\par
24 














Відношення $C$ між множинами $A$ та $B$ ($A, B\subset Z$) задане списком пар. Перевірити чи є відношення $C$ ін'єктивним.
\\\par
25 














Бінарне відношення R задане на множині $\overline{0,n-1}$ матрицею відношення з цілими елементами.  Перевірити, чи є відношення рефлексивним. Обов'язково зазначити, яке зображення матриці використовується в розв'язанні: двовимірний масив типу \textbf{ndarray} або список списків.
\\\par
26 














За розбиттям множини $A=\overline{0,n-1}$ та відображенням $f$, заданим на множині $A$, побудувати підрозбиття таке, що два елементи знаходяться в одному класі підрозбиття тоді і тільки тоді, коли відображення $f$ приймає на них однакове значення. Вважати, що $f$ задано функцією. Розбиття зображувати векторами (список або  \textbf{ndarray}), що задають належність елементів класам.
\\\par
27 














Бінарні відношення $R_1$, $R_2$ задано на множині $\overline{0,n-1}$ матрицями відношення з елементами логічного типу. Обчислити матрицю відношення $R_1\cup R_2$. Обов'язково зазначити, яке зображення матриці використовується в розв'язанні: двовимірний масив типу \textbf{ndarray} або список списків. Для всіх матриць має використовуватись однакове зображення.
\\\par
28 














Бінарне відношення на множині $\overline{0,n-1}$ задане множиною пар. Перевірити, чи є відношення рефлексивним.
\\\par
29 














Написати функцію, що перевіряє чи складаються дві послідовності з однакових елементів з точністю до кількості входжень кожного елемента в послідовність. Порядок елементів у послідовності несуттєвий.
\\\par
\newpage
30 














Бінарне відношення $R$ задане на множині $\overline{0,n-1}$ матрицею відношення з елементами логічного типу. Перевірити, чи є відношення антирефлексивним. Обов'язково зазначити, яке зображення матриці використовується в розв'язанні: двовимірний масив типу \textbf{ndarray} або список списків.
\\\par
31 














Бінарні відношення $R_1$, $R_2$ задано на множині $\overline{0,n-1}$ матрицями відношення з елементами логічного типу. Перевірити, чи правда, що  $R_1\subseteq R_2$. Обов'язково зазначити, яке зображення матриці використовується в розв'язанні: двовимірний масив типу \textbf{ndarray} або список списків.
\\\par
32 














Бінарне відношення R задане на множині $\overline{0,n-1}$ матрицею відношення з цілими елементами.  Перевірити, чи є відношення транзитивним. Обов'язково зазначити, яке зображення матриці використовується в розв'язанні: двовимірний масив типу \textbf{ndarray} або список списків.
\\\par
33 














Визначити, чи є одне розбиття множини $A=\overline{0,n-1}$\ підрозбиттям іншого. Розбиття задано векторами (список або  \textbf{ndarray}), що задають належність елементів класам.
\\\par
34 














Сімейство множин це множина, елементами якої є множини. Надалі будемо вважати, що сімейство множин задається як множина, елементами якої є незмінювані множини. Написати функцію, що перевіряє, чи є задане сімейство множин розбиттям заданої множини $A$.
\\\par
35 














За еквівалентністю на множині $\overline{0,n-1}$ побудувати її фактор-множину як множину, елементами якої є множини.
\\\par
36 














За розбиттям множини, що задане як множина множин, побудувати відповідну еквівалентність. Еквівалентність задати списком пар.
\\\par
37 














Визначити, чи є одне розбиття підрозбиттям іншого. Розбиття задано як словники, що задають належність елементів класам.
\\\par
38 














Функція $f$ типу $A\times B \rightarrow C$ задана словником. Розглядати "обернену" функцію $g$ з $C \times B$  в $P(A)$ (булеан множини $A$), $g\left(c, b\right)=\left\{a\middle|\ f\left(a,b\right)=c\right\}$. Написати функцію, що за функцією $f$ типу $Z\times Z \rightarrow Z$, знаходить "обернену" до неї функцію $g$ як функцію з $Z\times Z$ в $P(Z)$. Вважати, що функцію $f$ задано словником, результатом функції також має бути словник.
\\\par
39 














Розглядаються неспадні послідовності цілих чисел, що зберігаються в списках. Написати функцію, що перевіряє чи складаються дві послідовності з однакових елементів. Кількість входжень елемента в послідовність до уваги не брати. Наприклад, послідовності 1, 1, 3 та 1, 1, 1, 3, 3 складаються з однакових елементів. Обмеження: використовувати додаткові структури даних для збереження послідовностей або їх частин заборонено.
\\\par
40 














Для множин з повтореннями, заданими словниками, реалізувати теоретико-множинну операцію різниці. Розв’язання оформити у вигляді функції.
\\\par
41 














Рядок містить слова, відокремлені крапками та комами. Знайти довжину найдовшого слова рядка. Копій частини рядка довших 1 символу створювати не можна.
\\\par
42 














Бінарні відношення $R_1$, $R_2$ задано на множині $\overline{0,n-1}$ матрицями відношення з елементами логічного типу. Обчислити матрицю відношення $R_1\cap R_2$. Обов'язково зазначити, яке зображення матриці використовується в розв'язанні: двовимірний масив типу \textbf{ndarray} або список списків. Для всіх матриць має використовуватись однакове зображення.
\\\par
\newpage
43 














Розглядаються неспадні послідовності цілих чисел, що зберігаються в списках. Написати функцію, що утворює нову неспадну послідовність, яка є об'єднанням двох своїх аргументів. Якщо елемент в першу послідовність входить $k_1$ раз, а в другу --- $k_2$ разів, то в результат він має входити $max\left\{k_1,k_2\right\}$ разів. Наприклад, за послідовностей 1, 3, 3, 5 та 2, 3, 5, 8 результатом має бути 1, 2, 3, 3, 5, 8.
\\\par
44 














Для множин з повтореннями, заданими словниками, реалізувати теоретико-множинну операцію перевірки строгого включення.
\\\par
45 














Бінарні відношення $R_1$, $R_2$ задано на множині $\overline{0,n-1}$ матрицями відношення з цілими елементами. Обчислити матрицю відношення $R_1\cap R_2$. Обов'язково зазначити, яке зображення матриці використовується в розв'язанні: двовимірний масив типу \textbf{ndarray} або список списків. Для всіх матриць має використовуватись однакове зображення.
\\\par
46 














Бінарне відношення $R$ задане на множині $\overline{0,n-1}$ матрицею відношення. Побудувати обернене відношення. Матрицю задано як список списків.
\\\par
47 














Бінарне відношення на множині $\overline{0,n-1}$ задане множиною пар. Побудувати обернене відношення.
\\\par
48 














За еквівалентністю на множині $A=\overline{0,n-1}$ та елементом множини $A$ обчислити його клас еквівалентності.
\\\par
49 














Бінарні відношення $R_1$, $R_2$ задано на множині $\overline{0,n-1}$ матрицями відношення з елементами логічного типу. Обчислити матрицю відношення $R_1\setminus R_2$. Обов'язково зазначити, яке зображення матриці використовується в розв'язанні: двовимірний масив типу \textbf{ndarray} або список списків. Для всіх матриць має використовуватись однакове зображення.
\\\par
50 














Бінарні відношення $R_1$, $R_2$ задано на множині $\overline{0,n-1}$ матрицями відношення з цілими елементами. Перевірити, чи правда, що  $R_1\subseteq R_2$. Обов'язково зазначити, яке зображення матриці використовується в розв'язанні: двовимірний масив типу \textbf{ndarray} або список списків.
\\\par
51 














За еквівалентністю на множині $\overline{0,n-1}$ побудувати її фактор-множину як вектор (список або  \textbf{ndarray}), що задає належність елементів класам.
\\\par
52 














Сімейство множин це множина, елементами якої є множини. Надалі будемо вважати, що сімейство множин задається як множина, елементами якої є незмінювані множини. Написати функцію, що за сімейством множин перевіряє, чи правда, що елементи сімейства попарно не перетинаються.
\\\par
53 














Написати функцію, що за послідовністю чисел будує список її елементів, що зустрічаються в ній рівно один раз. Числа, що мають різний тип, вважаємо різними.
\\\par
54 














За послідовністю цілих чисел s визначити яких значень в ній більше: парних чи непарних. Кожне значення враховується рівно один раз. Якщо більше парних, повернути 2, якщо більше непарних повернути 1, якщо кількості однакові, повернути 0. Розв’язок оформити у вигляді функції.
\\\par
55 














Для функції $f$ типу $A\rightarrow B$ обернена функція, узагалі-то кажучи, функцією з $B\rightarrow A$ може не бути. Тим не менш можна розглядати "обернену" функцію $g$ з $B$ в $P(A)$ (булеан множини $A$), $g\left(b\right)=f^{-1}\left(b\right)$ ($f^{-1}\left(b\right)$ -- повний прообраз елемента $b$). Написати функцію, що за функцією $f$  типу $Z\rightarrow Z$, знаходить "обернену" до неї функцію $g$ як функцію з $Z$ в $P(Z)$. Вважати, що функцію $f$ задано словником, результатом функції також має бути словник.
\\\par
56 














Написати функцію, що перевіряє чи є послідовність дійсних чисел монотонною.
\\\par
\newpage
57 














Рядок містить слова, відокремлені крапками та комами. Знайти останнє найдовше слово рядка. Копії частини рядка довші 1 символу можна створювати тільки в інструкції \textbf{return}.
\\\par
58 














Бінарні відношення $R_1$, $R_2$ задано на множині $\overline{0,n-1}$ матрицями відношення з цілими елементами. Обчислити матрицю відношення $R_1\setminus R_2$. Обов'язково зазначити, яке зображення матриці використовується в розв'язанні: двовимірний масив типу \textbf{ndarray} або список списків. Для всіх матриць має використовуватись однакове зображення.
\\\par
59 














Написати функцію, що за словником обчислює скільки разів зустрічається в ньому задане значення.
\\\par
60 














Побудувати функцію, що зображення розбиття множини\\ $\overline{0,n-1}$  у вигляді множини множин перетворює на зображення розбиття у вигляді вектору (список або  \textbf{ndarray}), що задає належність елементів класам.
\\\par
61 














Бінарне відношення $R$ задане на множині $\overline{0,n-1}$ матрицею відношення з елементами логічного типу. Перевірити, чи є відношення антисиметричним. Обов'язково зазначити, яке зображення матриці використовується в розв'язанні: двовимірний масив типу \textbf{ndarray} або список списків.
\\\par
62 














За розбиттям множини $A=\overline{0,n-1}$, що задане як вектор (список або  \textbf{ndarray}), що задає належність елементів класам, побудувати відповідну еквівалентність. Еквівалентність задати матрицею відношення.
\\\par
63 














Рядок містить слова, відокремлені крапками та комами. Знайти кількість слів, довжина яких більша 3. Копій частини рядка довших 1 символу створювати не можна.
\\\par
64 














Написати функцію, що повертає кількість різних дійсних значень, що містяться в послідовності s.
\\\par
65 














Бінарне відношення $R$ задане на множині $\overline{0,n-1}$ матрицею відношення з елементами логічного типу. Перевірити, чи є відношення транзитивним. Обов'язково зазначити, яке зображення матриці використовується в розв'язанні: двовимірний масив типу \textbf{ndarray} або список списків.
\\\par
66 














Для множин з повтореннями, заданими словниками, реалізувати теоретико-множинну операцію злиття (елемент входить в результат суму входжень в аргументи). Розв’язання оформити у вигляді функції.
\\\par
67 














Бінарне відношення на множині $\overline{0,n-1}$ задане множиною пар. Перевірити, чи є відношення антирефлексивним.
\\\par
68 














Бінарні відношення $R_1$, $R_2$ задано на множині $\overline{0,n-1}$ матрицями відношення з цілими елементами. Обчислити композицію $R_1\circ R_2$ відношень. Обов'язково зазначити, яке зображення матриці використовується в розв'язанні: двовимірний масив типу \textbf{ndarray} або список списків.
\\\par
69 














Написати функцію, що за послідовністю цілих знаходить всі елементи, що зустрічаються в ній рівно 7 разів.
\\\par
70 














Словники a1 та a2 містять канонічні розклади натуральних чисел n1 та n2 на прості множники. Написати функцію, що знаходить канонічний розклад найменшого спільного кратного цих чисел.
\\\par
71 














Бінарні відношення $R_1$, $R_2$ задано на множині $\overline{0,n-1}$ матрицями відношення з цілими елементами. Перевірити, чи правда, що $R_1\subset R_2$. Обов'язково зазначити, яке зображення матриці використовується в розв'язанні: двовимірний масив типу \textbf{ndarray} або список списків.
\\\par
\newpage
72 














Список містить цілі числа. Вилучити зі списку усі повторення (кожне значення має залишитися в одному екземплярі). Використання додаткових структур даних заборонено. Для кожного елемента в списку має залишитись його останнє входження і при цьому порядок елементів змінювати не можна.
\\\par
73 














Побудувати функцію, що зображення розбиття множини у вигляді словника, що задає належність елементів класам, перетворює на зображення у вигляді множини множин.
\\\par
74 














Сімейство множин це множина, елементами якої є множини. Надалі будемо вважати, що сімейство множин задається як множина, елементами якої є незмінювані множини. Написати функцію, що перевіряє, чи є задане сімейство множин покриттям заданої множини $A$.
\\\par
75 














Сімейство множин це множина, елементами якої є множини. Надалі будемо вважати, що сімейство множин задається як множина, елементами якої є незмінювані множини. Написати функцію, що обчислює перетин сімейства множин. За спроби обчислити перетин порожнього сімейства має генеруватись виняток.
\\\par
76 














Розглядаються неспадні послідовності цілих чисел, що зберігаються в списках. Написати функцію, що утворює нову неспадну послідовність, яка є перетином двох своїх аргументів. Якщо елемент в першу послідовність входить $k_1$ раз, а в другу --- $k_2$ разів, то в результат він має входити $min\left\{k_1,k_2\right\}$ разів. Наприклад, за послідовностей 1, 3, 3, 5 та 2, 3, 3, 3, 5, 8 результатом має бути 3, 3, 5.
\\\par
77 














Розглядаються неспадні послідовності цілих чисел, що зберігаються в списках. Написати функцію, що утворює нову неспадну послідовність, яка є злиттям двох своїх аргументів. Якщо елемент в першу послідовність входить $k_1$ раз, а в другу --- $k_2$ разів, то в результат він має входити $k_1+k_2$ разів. Наприклад, за послідовностей 1, 3, 3, 5 та 2, 3, 5, 8 результатом має бути 1, 2, 3, 3, 3, 5, 5, 8.
\\\par
78 














Бінарні відношення $R_1$, $R_2$ задано на множині $\overline{0,n-1}$ матрицями відношення з цілими елементами. Обчислити матрицю відношення $R_1\cup R_2$. Обов'язково зазначити, яке зображення матриці використовується в розв'язанні: двовимірний масив типу \textbf{ndarray} або список списків. Для всіх матриць має використовуватись однакове зображення.
\\\par
79 














Словники a1 та a2 містять канонічні розклади натуральних чисел n1 та n2 на прості множники. Написати функцію, що знаходить найменше спільне кратне цих чисел.
\\\par
80 














Бінарне відношення R задане на множині $\overline{0,n-1}$ матрицею відношення з цілими елементами.  Перевірити, чи є відношення антисиметричним. Обов'язково зазначити, яке зображення матриці використовується в розв'язанні: двовимірний масив типу \textbf{ndarray} або список списків.
\\\par
81 














Бінарне відношення $R$ задане на множині $\overline{0,n-1}$ матрицею відношення з елементами логічного типу. Перевірити, чи є відношення рефлексивним. Обов'язково зазначити, яке зображення матриці використовується в розв'язанні: двовимірний масив типу \textbf{ndarray} або список списків.
\\\par
82 














Словники a1 та a2 містять канонічні розклади натуральних чисел n1 та n2 на прості множники. Написати функцію, що знаходить найбільший спільний дільник цих чисел.
\\\par
83 














Відношення $C$ між множинами $A$ та $B$ ($A, B\subset Z$) задане списком пар. Перевірити чи є відношення $C$ функціональним.
\\\par
84 














Словники a1 та a2 містять канонічні розклади натуральних чисел n1 та n2 на прості множники. Написати функцію, що знаходить канонічний розклад найбільшого спільного дільника цих чисел.
\\\par
\newpage
85 














Сімейство множин це множина, елементами якої є множини. Надалі будемо вважати, що сімейство множин задається як множина, елементами якої є незмінювані множини. Написати функцію, що обчислює об'єднання сімейства множин. Об'єднання порожнього сімейства вважати порожньою множиною.
\\\par
86 














Бінарне відношення $R$ задане на множині $\overline{0,n-1}$ матрицею відношення з елементами логічного типу. Перевірити, чи є відношення симетричним. Обов'язково зазначити, яке зображення матриці використовується в розв'язанні: двовимірний масив типу \textbf{ndarray} або список списків.
\\\par
87 














За розбиттям множини $A=\overline{0,n-1}$, що задане як вектор (список або  \textbf{ndarray}), що задає належність елементів класам, побудувати відповідну еквівалентність. Еквівалентність задати списком пар.
\\\par
88 














Список містить числа. Вилучити зі списку усі повторення (кожне значення має залишитися в одному екземплярі). Для кожного числа має залишатись його перше входження в список. Числа, що мають різний тип, вважаємо різними.
\\\par
89 














За розбиттям множини, що задане як множина множин, побудувати відповідну еквівалентність. Еквівалентність задати матрицею відношення.
\\\par
90 














Список містить цілі числа. Вилучити зі списку усі повторення (кожне значення має залишитися в одному екземплярі). Використання додаткових структур даних заборонено. Для кожного елемента в списку має залишитись його перше входження і при цьому порядок елементів змінювати не можна.
\\\par
91 














Для множин з повтореннями, заданими словниками, реалізувати теоретико-множинну операцію перевірки включення. Розв'яз\-ан\-ня оформити у вигляді функції.
\\\par
92 














Словник a містить канонічний розклад натурального числа n на прості множники. Написати функцію, що за словником a обчислює число n.
\\\par
93 














Функція $f$ між множинами $A$ та $B$ (елементи множини $A$ є такими, що хешуються) задана словником, в якому значеннями ключів є значення функції. Написати функцію, що зображує $f$ у вигляді списку пар.
\\\par
94 














Побудувати функцію, що зображення розбиття множини\\ $\overline{0,n-1}$  у вигляді вектору (список або  \textbf{ndarray}), що задає належність елементів класам, перетворює на зображення у вигляді множини множин.
\\\par
95 














За функцією $f$, заданою словником, та елементом $b$ знайти значення $f^{-1}(b)$. Розв’язання оформити у вигляді функції, що приймає $f$ та $b$ як аргументи.
\\\par
96 














Бінарне відношення на множині $\overline{0,n-1}$ задане множиною пар. Перевірити, чи є відношення транзитивним.
\\\par
97 














Бінарне відношення R задане на множині $\overline{0,n-1}$ матрицею відношення з цілими елементами.  Перевірити, чи є відношення транзитивним. Обов'язково зазначити, яке зображення матриці використовується в розв'язанні: двовимірний масив типу \textbf{ndarray} або список списків.
\\\par
98 














Бінарні відношення $R_1$, $R_2$ задано на множині $\overline{0,n-1}$ матрицями відношення з цілими елементами. Обчислити матрицю відношення $R_1\setminus R_2$. Обов'язково зазначити, яке зображення матриці використовується в розв'язанні: двовимірний масив типу \textbf{ndarray} або список списків. Для всіх матриць має використовуватись однакове зображення.
\\\par
\newpage
99 














Бінарне відношення $R$ задане на множині $\overline{0,n-1}$ матрицею відношення з елементами логічного типу. Перевірити, чи є відношення симетричним. Обов'язково зазначити, яке зображення матриці використовується в розв'язанні: двовимірний масив типу \textbf{ndarray} або список списків.
\\\par
100 














За еквівалентністю на множині $\overline{0,n-1}$ побудувати її фактор-множину як множину, елементами якої є множини.
\\\par
101 














Бінарне відношення R задане на множині $\overline{0,n-1}$ матрицею відношення з цілими елементами.  Перевірити, чи є відношення антисиметричним. Обов'язково зазначити, яке зображення матриці використовується в розв'язанні: двовимірний масив типу \textbf{ndarray} або список списків.
\\\par
102 














Бінарні відношення $R_1$, $R_2$ задано на множині $\overline{0,n-1}$ матрицями відношення з цілими елементами. Перевірити, чи правда, що $R_1\subset R_2$. Обов'язково зазначити, яке зображення матриці використовується в розв'язанні: двовимірний масив типу \textbf{ndarray} або список списків.
\\\par
103 














Бінарні відношення $R_1$, $R_2$ задано на множині $\overline{0,n-1}$ матрицями відношення з цілими елементами. Обчислити композицію $R_1\circ R_2$ відношень. Обов'язково зазначити, яке зображення матриці використовується в розв'язанні: двовимірний масив типу \textbf{ndarray} або список списків.
\\\par
104 














Сімейство множин це множина, елементами якої є множини. Надалі будемо вважати, що сімейство множин задається як множина, елементами якої є незмінювані множини. Написати функцію, що обчислює об'єднання сімейства множин. Об'єднання порожнього сімейства вважати порожньою множиною.
\\\par
105 














Функція $f$ між множинами $A$ та $B$ (елементи множини $A$ є такими, що хешуються) задана словником, в якому значеннями ключів є значення функції. Написати функцію, що зображує $f$ у вигляді списку пар.
\\\par
106 














За функціями $f$ та $g$, заданими словниками знайти їх композицію. Розв’язання оформити у вигляді функції, що приймає словники $f$ та $g$ як аргументи та результат повертає у вигляді словника.
\\\par
107 














Список містить цілі числа. Вилучити зі списку усі повторення (кожне значення має залишитися в одному екземплярі). Для кожного числа має залишатись його перше входження в список.
\\\par
108 














Розглядаються неспадні послідовності цілих чисел, що зберігаються в списках. Написати функцію, що перевіряє чи складаються дві послідовності з однакових елементів. Кількість входжень елемента в послідовність до уваги не брати. Наприклад, послідовності 1, 1, 3 та 1, 1, 1, 3, 3 складаються з однакових елементів. Обмеження: використовувати додаткові структури даних для збереження послідовностей або їх частин заборонено.
\\\par
109 














Розглядаються неспадні послідовності цілих чисел, що зберігаються в списках. Написати функцію, що утворює нову неспадну послідовність, яка є перетином двох своїх аргументів. Якщо елемент в першу послідовність входить $k_1$ раз, а в другу --- $k_2$ разів, то в результат він має входити $min\left\{k_1,k_2\right\}$ разів. Наприклад, за послідовностей 1, 3, 3, 5 та 2, 3, 3, 3, 5, 8 результатом має бути 3, 3, 5.
\\\par
110 














Бінарні відношення $R_1$, $R_2$ задано на множині $\overline{0,n-1}$ матрицями відношення з елементами логічного типу. Перевірити, чи правда, що  $R_1\subseteq R_2$. Обов'язково зазначити, яке зображення матриці використовується в розв'язанні: двовимірний масив типу \textbf{ndarray} або список списків.
\\\par
111 














Сімейство множин це множина, елементами якої є множини. Надалі будемо вважати, що сімейство множин задається як множина, елементами якої є незмінювані множини. Написати функцію, що за сімейством множин перевіряє, чи правда, що елементи сімейства попарно не перетинаються.
\\\par
\newpage
112 














Визначити, чи є одне розбиття множини $A=\overline{0,n-1}$\ підрозбиттям іншого. Розбиття задано векторами (список або  \textbf{ndarray}), що задають належність елементів класам.
\\\par
113 














За розбиттям множини, що задане як множина множин, побудувати відповідну еквівалентність. Еквівалентність задати списком пар.
\\\par
114 














Для множин з повтореннями, заданими словниками, реалізувати теоретико-множинну операцію різниці. Розв’язання оформити у вигляді функції.
\\\par
115 














Бінарне відношення $R$ задане на множині $\overline{0,n-1}$ матрицею відношення з елементами логічного типу. Перевірити, чи є відношення рефлексивним. Обов'язково зазначити, яке зображення матриці використовується в розв'язанні: двовимірний масив типу \textbf{ndarray} або список списків.
\\\par
116 














Бінарне відношення на множині $\overline{0,n-1}$ задане множиною пар. Перевірити, чи є відношення симетричним.
\\\par
117 














Рядок містить слова, відокремлені крапками та комами. Знайти кількість слів, довжина яких дорівнює 13. Копій частини рядка довших 1 символу створювати не можна.
\\\par
118 














Бінарне відношення на множині $\overline{0,n-1}$ задане множиною пар. Побудувати обернене відношення.
\\\par
119 














Бінарне відношення $R$ задане на множині $\overline{0,n-1}$ матрицею відношення з елементами логічного типу. Перевірити, чи є відношення транзитивним. Обов'язково зазначити, яке зображення матриці використовується в розв'язанні: двовимірний масив типу \textbf{ndarray} або список списків.
\\\par
120 














Бінарні відношення $R_1$, $R_2$ задано на множині $\overline{0,n-1}$ матрицями відношення з цілими елементами. Перевірити, чи правда, що  $R_1\subseteq R_2$. Обов'язково зазначити, яке зображення матриці використовується в розв'язанні: двовимірний масив типу \textbf{ndarray} або список списків.
\\\par
121 














За розбиттям множини $A=\overline{0,n-1}$ та відображенням\\ $f{:}A\rightarrow{bool}$  побудувати підрозбиття таке, що два елементи знаходяться в одному класі підрозбиття тоді і тільки тоді, коли відображення $f$ приймає на них однакове значення. Вважати, що $f$ задано функцією. Для розбиттів використовувати зображення у вигляді множини множин.
\\\par
122 














Побудувати функцію, що зображення розбиття множини\\ $\overline{0,n-1}$  у вигляді вектору (список або  \textbf{ndarray}), що задає належність елементів класам, перетворює на зображення у вигляді множини множин.
\\\par
123 














Рядок містить слова, відокремлені крапками та комами. Знайти кількість слів, що містять кожну з літер s, e, l, f. Копій частини рядка довших 1 символу створювати не можна.
\\\par
124 














Бінарне відношення R задане на множині $\overline{0,n-1}$ матрицею відношення з цілими елементами.  Перевірити, чи є відношення рефлексивним. Обов'язково зазначити, яке зображення матриці використовується в розв'язанні: двовимірний масив типу \textbf{ndarray} або список списків.
\\\par
125 














Бінарні відношення $R_1$, $R_2$ задано на множині $\overline{0,n-1}$ матрицями відношення з елементами логічного типу. Обчислити матрицю відношення $R_1\cup R_2$. Обов'язково зазначити, яке зображення матриці використовується в розв'язанні: двовимірний масив типу \textbf{ndarray} або список списків. Для всіх матриць має використовуватись однакове зображення.
\\\par
\newpage
126 














Відповідність $C$ між множинами $A$ та $B$ ($A\subset Z$) задана списком пар. Написати функцію, що зображує $C$ у вигляді словника, де значеннями ключів є їх повні образи.
\\\par

127 














За функцією $f$, заданою словником, та елементом $b$ знайти значення $f^{-1}(b)$. Розв’язання оформити у вигляді функції, що приймає $f$ та $b$ як аргументи.
\\\par
128 














Рядок містить слова, відокремлені крапками та комами. Знайти кількість слів, довжина яких більша 3. Копій частини рядка довших 1 символу створювати не можна.
\\\par
129 














Словники a1 та a2 містять канонічні розклади натуральних чисел n1 та n2 на прості множники. Написати функцію, що знаходить найбільший спільний дільник цих чисел.
\\\par
130 














За розбиттям множини, що задане як множина множин, побудувати відповідну еквівалентність. Еквівалентність задати матрицею відношення.
\\\par
131 














Бінарне відношення на множині $\overline{0,n-1}$ задане множиною пар. Перевірити, чи є відношення антирефлексивним.
\\\par
132 














Написати функцію, що за послідовністю цілих знаходить 5 елементів, що зустрічаються в ній найчастіше. Якщо п’яте місце розділяють кілька значень, додавати до результату всіх їх.
\\\par
133 














Написати функцію, що за послідовністю цілих знаходить елементи, що зустрічаються в ній рідше за інші (але зустрічаються).
\\\par
134 














Бінарні відношення $R_1$, $R_2$ задано на множині $\overline{0,n-1}$ матрицями відношення з елементами логічного типу. Перевірити, чи правда, що $R_1\subset R_2$. Обов'язково зазначити, яке зображення матриці використовується в розв'язанні: двовимірний масив типу \textbf{ndarray} або список списків.
\\\par
135 














Бінарні відношення $R_1$, $R_2$ задано на множині $\overline{0,n-1}$ матрицями відношення з цілими елементами. Обчислити матрицю відношення $R_1\cap R_2$. Обов'язково зазначити, яке зображення матриці використовується в розв'язанні: двовимірний масив типу \textbf{ndarray} або список списків. Для всіх матриць має використовуватись однакове зображення.
\\\par
136 














Розглядаються неспадні послідовності цілих чисел, що зберігаються в списках. Написати функцію, що утворює нову неспадну послідовність, яка є злиттям двох своїх аргументів. Якщо елемент в першу послідовність входить $k_1$ раз, а в другу --- $k_2$ разів, то в результат він має входити $k_1+k_2$ разів. Наприклад, за послідовностей 1, 3, 3, 5 та 2, 3, 5, 8 результатом має бути 1, 2, 3, 3, 3, 5, 5, 8.
\\\par
137 














Сімейство множин це множина, елементами якої є множини. Надалі будемо вважати, що сімейство множин задається як множина, елементами якої є незмінювані множини. Написати функцію, що перевіряє, чи є задане сімейство множин покриттям заданої множини $A$.
\\\par
138 














Список містить цілі числа. Вилучити зі списку усі повторення (кожне значення має залишитися в одному екземплярі). Використання додаткових структур даних заборонено. Для кожного елемента в списку має залишитись його перше входження і при цьому порядок елементів змінювати не можна.
\\\par
139 














Функція $f$ між множинами $A$ та $B$ (елементи множини $A$ є такими, що хешуються) задана списком пар. Написати функцію, що зображує $f$ у вигляді словника, де значеннями ключів є значення функції.
\\\par
\newpage
140 














Список містить цілі числа. Вилучити зі списку усі повторення (кожне значення має залишитися в одному екземплярі). Використання додаткових структур даних заборонено. Для кожного елемента в списку має залишитись його останнє входження і при цьому порядок елементів змінювати не можна.
\\\par
141 














За розбиттям множини $A=\overline{0,n-1}$ та відображенням $f$, заданим на множині $A$, побудувати підрозбиття таке, що два елементи знаходяться в одному класі підрозбиття тоді і тільки тоді, коли відображення $f$ приймає на них однакове значення. Вважати, що $f$ задано функцією. Розбиття зображувати векторами (список або  \textbf{ndarray}), що задають належність елементів класам.
\\\par
142 














За еквівалентністю на множині $A=\overline{0,n-1}$ та елементом множини $A$ обчислити його клас еквівалентності.
\\\par
143 














Бінарне відношення R задане на множині $\overline{0,n-1}$ матрицею відношення з цілими елементами.  Перевірити, чи є відношення антирефлексивним. Обов'язково зазначити, яке зображення матриці використовується в розв'язанні: двовимірний масив типу \textbf{ndarray} або список списків.
\\\par
144 














Написати функцію, що за послідовністю чисел будує список її елементів, що зустрічаються в ній рівно один раз. Числа, що мають різний тип, вважаємо різними.
\\\par
145 














Бінарне відношення $R$ задане на множині $\overline{0,n-1}$ матрицею відношення. Побудувати обернене відношення. Матрицю задано як список списків.
\\\par
146 














Відношення $C$ між множинами $A$ та $B$ ($A, B\subset Z$) задане словником, в якому ключами є елементи множини $A$, а значенням кожного ключа $a$ --- його повний образ $C\left(a\right)=\left\{b|\left(a,b\right)\in C\right\}$. Перевірити чи є відношення $C$ ін'єктивним.
\\\par
147 














Сімейство множин це множина, елементами якої є множини. Надалі будемо вважати, що сімейство множин задається як множина, елементами якої є незмінювані множини. Написати функцію, що перевіряє, чи є задане сімейство множин розбиттям заданої множини $A$.
\\\par
148 














Побудувати функцію, що зображення розбиття множини у вигляді словника, що задає належність елементів класам, перетворює на зображення у вигляді множини множин.
\\\par
149 














Бінарні відношення $R_1$, $R_2$ задано на множині $\overline{0,n-1}$ матрицями відношення з цілими елементами. Обчислити матрицю відношення $R_1\cup R_2$. Обов'язково зазначити, яке зображення матриці використовується в розв'язанні: двовимірний масив типу \textbf{ndarray} або список списків. Для всіх матриць має використовуватись однакове зображення.
\\\par
150 














Написати функцію, що стискає послідовність цілих, зображену списком: нулі в кінці та на початку вилучаються, у середині послідовності кожна послідовність нулів замінюється на один нуль. Наприклад, з послідовності 0,0,1,2,0,0,0,0,3,0,5,0,0 має стати 1,2,0,3,0,5. Нові об'єкти, що зберігають послідовності утворювати не можна.
\\\par
151 














Написати функцію, що за послідовністю цілих знаходить всі елементи, що зустрічаються в ній рівно 7 разів.
\\\par
152 














Визначити, чи є одне розбиття підрозбиттям іншого. Розбиття задано як словники, що задають належність елементів класам.
\\\par
153 














За розбиттям множини $A=\overline{0,n-1}$, що задане як вектор (список або  \textbf{ndarray}), що задає належність елементів класам, побудувати відповідну еквівалентність. Еквівалентність задати матрицею відношення.
\\\par
\newpage
154 














Написати функцію, що повертає кількість різних дійсних значень, що містяться в послідовності s.
\\\par
155 














Словник a містить канонічний розклад натурального числа n на прості множники. Написати функцію, що за словником a обчислює число n.
\\\par
156 














Для множин з повтореннями, заданими словниками, реалізувати теоретико-множинну операцію перевірки включення. Розв'яз\-ан\-ня оформити у вигляді функції.
\\\par
157 














За розбиттям множини $A=\overline{0,n-1}$, що задане як вектор (список або  \textbf{ndarray}), що задає належність елементів класам, побудувати відповідну еквівалентність. Еквівалентність задати списком пар.
\\\par
158 














Список містить числа. Вилучити зі списку усі повторення (кожне значення має залишитися в одному екземплярі). Для кожного числа має залишатись його перше входження в список. Числа, що мають різний тип, вважаємо різними.
\\\par
159 














Для множин з повтореннями, заданими словниками, реалізувати теоретико-множинну операцію об’єднання (елемент входить у результат максимум входжень в аргументи). Розв’язання оформити у вигляді функції.
\\\par
160 














Функція $f$ типу $A\times B \rightarrow C$ задана словником. Розглядати "обернену" функцію $g$ з $C \times B$  в $P(A)$ (булеан множини $A$), $g\left(c, b\right)=\left\{a\middle|\ f\left(a,b\right)=c\right\}$. Написати функцію, що за функцією $f$ типу $Z\times Z \rightarrow Z$, знаходить "обернену" до неї функцію $g$ як функцію з $Z\times Z$ в $P(Z)$. Вважати, що функцію $f$ задано словником, результатом функції також має бути словник.
\\\par
161 














За еквівалентністю на множині $\overline{0,n-1}$ побудувати її фактор-множину як вектор (список або  \textbf{ndarray}), що задає належність елементів класам.
\\\par
162 














Для множин з повтореннями, заданими словниками, реалізувати теоретико-множинну операцію перетину (елемент входить в результат мінімум входжень в аргументи). Розв’язання оформити у вигляді функції.
\\\par
163 














Бінарні відношення $R_1$, $R_2$ задано на множині $\overline{0,n-1}$ матрицями відношення з елементами логічного типу. Обчислити матрицю відношення $R_1\setminus R_2$. Обов'язково зазначити, яке зображення матриці використовується в розв'язанні: двовимірний масив типу \textbf{ndarray} або список списків. Для всіх матриць має використовуватись однакове зображення.
\\\par
164 














Побудувати функцію, що зображення розбиття множини у вигляді множини множин перетворює на зображення розбиття у вигляді словника, що задає належність елементів класам.
\\\par
165 














Бінарне відношення на множині $\overline{0,n-1}$ задане множиною пар. Перевірити, чи є відношення антисиметричним.
\\\par
166 














Відношення $C$ між множинами $A$ та $B$ ($A, B\subset Z$) задане списком пар. Перевірити чи є відношення $C$ ін'єктивним.
\\\par
167 














Бінарне відношення на множині $\overline{0,n-1}$ задане множиною пар. Перевірити, чи є відношення транзитивним.
\\\par
168 














Словники a1 та a2 містять канонічні розклади натуральних чисел n1 та n2 на прості множники. Написати функцію, що знаходить канонічний розклад найменшого спільного кратного цих чисел.
\\\par
\newpage
169 














Бінарні відношення $R_1$, $R_2$ задано на множині $\overline{0,n-1}$ матрицями відношення з елементами логічного типу. Обчислити композицію $R_1\circ R_2$ відношень. Обов'язково зазначити, яке зображення матриці використовується в розв'язанні: двовимірний масив типу \textbf{ndarray} або список списків.
\\\par
170 














Визначити, чи є одне розбиття підрозбиттям іншого. Розбиття задані як множини множин.
\\\par
